\chapter{Bence-Jones Protein (Urine Immunofixation)}

\section*{What Is This Test?}
Bence-Jones protein is monoclonal free immunoglobulin light chain excreted in urine. Urine protein electrophoresis and immunofixation can detect and characterize kappa or lambda light chains produced by plasma-cell disorders.

\section*{Alternative Names}
Urine Immunofixation; Urine Free Light Chains; Urine Bence-Jones Protein; UPEP/Urine IFE.

\section*{Why This Test Is Ordered}
Testing may be used in suspected light-chain multiple myeloma, AL amyloidosis, light-chain deposition disease, or cast nephropathy, particularly when serum protein electrophoresis is negative or when significant proteinuria needs characterization. Serum free light chains are generally preferred for initial screening, with urine testing used selectively.

\section*{How This Test Is Performed}
The laboratory analyzes a timed urine collection or an appropriately collected urine sample by electrophoresis and immunofixation. The pattern identifies monoclonal kappa or lambda light chains and may quantify total urinary protein excretion.

\section*{Preparation, Risks, and Limitations}
Follow collection instructions carefully and do not discard specimens from a timed collection. The test is noninvasive. Kidney function, urine concentration, incomplete collection, and low-level secretion can affect sensitivity. Results should be interpreted with serum electrophoresis, serum immunofixation, and serum free light chains.

\section*{Understanding Results}
A monoclonal light-chain band supports a plasma-cell or B-cell clone but does not by itself establish multiple myeloma. The amount, type, serum findings, kidney function, blood counts, calcium, and clinical features determine the need for further evaluation.

\section*{References}
International Myeloma Working Group guidance on laboratory evaluation of monoclonal gammopathies.

\clearpage
\section*{Understanding Results and Follow-up}
The laboratory report should identify the specimen, method, units, reference interval or decision limit, and whether the result is positive, negative, abnormal, or inconclusive. Reference intervals vary with age, sex, pregnancy, specimen type, assay, and laboratory; a result outside a range is not automatically a diagnosis, and a result within a range does not always exclude disease. Discuss unexpected results with the ordering clinician, who can decide whether repeat or confirmatory testing, treatment, or referral is appropriate. Do not change medicines or delay urgent care based on an isolated result.


\section*{Regulatory Status and Authorization Date}
Regulatory status applies to the specific assay, instrument, manufacturer, and intended use rather than to the test name alone. No single approval date can be assigned to this test category. For a commercial assay, verify the current product in the relevant regulator's database: the U.S. Food and Drug Administration (FDA) clearance, approval, or authorization; India's Central Drugs Standard Control Organisation (CDSCO) licence or permission; the European Union In Vitro Diagnostic Regulation (EU IVDR) CE certificate issued through the applicable conformity-assessment route; or the United Kingdom Medicines and Healthcare products Regulatory Agency (MHRA) registration and UKCA/CE status. Laboratory-developed tests may not have an individual FDA, CDSCO, EU IVDR, or MHRA approval date.
\clearpage
